% Font
%\usepackage[T1]{fontenc}
%\renewcommand*\familydefault{\sfdefault} %% Only if the base font of the document is to be sans serif

\usepackage[colorlinks=true,urlcolor=magenta,bookmarksopen=true]{hyperref} 	% Záložky v~souboru
\usepackage[a-2u]{pdfx}
\usepackage[czech]{babel}
% \usepackage[defaultfam,tabular,lining]{montserrat}
\usepackage{lmodern}
% \usepackage[utf8]{inputenc}
\usepackage[T1]{fontenc}
\usepackage{textcomp}
\usepackage[linesnumbered,ruled,vlined]{algorithm2e}                % Pseudokód algoritmů
\usepackage{ifthen}

%\usepackage{times}						% Font (~ Times New Roman)

\usepackage{pifont}
\usepackage{tcolorbox}
\usepackage{multicol}
\usepackage{wrapfig}
\usepackage{caption}
\usepackage{subcaption}
\usepackage{geometry} 					% Odsazení
\usepackage{csquotes}					% české uvozovky
\usepackage{enumitem}					% enumerate environment
\usepackage{lipsum}						% Lorem ipsum
\usepackage{amsmath}
\usepackage{amssymb}
\usepackage{amsthm}						% Věta, lemma, důsledek, důkaz
\usepackage{tikz}						% Schémata automatů
\IfFileExists{fontawesome5.sty}{%
    \usepackage{awesomebox}%            % infoboxíky
}{%
    % Jednoduchá náhrada pro minimální instalace TeX Live.
    \newcommand{\faInfoCircle}{\ensuremath{\mathrm{i}}}
    \newcommand{\awesomebox}[5][]{\fbox{\parbox{0.9\linewidth}{##5}}}
}

\usepackage{pgfplots}
\pgfplotsset{compat=1.15}
\usepackage{mathrsfs}
\usetikzlibrary{arrows}

% Font
\IfFileExists{newtxtext.sty}{%
    \IfFileExists{mtpro2.sty}{%
        \usepackage{newtxtext}
        \usepackage[lite,subscriptcorrection,slantedGreek,nofontinfo,amsbb,eucal]{mtpro2}
    }{}%
}{}% V minimální instalaci zůstane výše načtený Latin Modern.

% Mezery v obsahu odpovídají sazbě projektu AM-skripta. Balíček
% tocloft ponechává mezi kapitolami větší mezeru, zatímco položky
% sekcí a podsekcí sází kompaktně pod sebe.
\usepackage{tocloft}
\ifcsname cftbeforechapskip\endcsname
    \setlength{\cftbeforechapskip}{1em plus 1pt}
\fi
\setlength{\cftbeforesecskip}{0pt plus 0.2pt}
\setlength{\cftbeforesubsecskip}{0pt plus 0.2pt}

% Indentation
\newcommand{\hmargin}{18mm}
\newcommand{\vmargin}{25mm}

% Image sizing
\newcommand{\imgwidthpercent}{.75}
\newcommand{\imagespacing}{3em}
\newcommand{\cubewidth}{.4\textwidth}
\newcommand{\dotswidth}{.5\textwidth}
\newcommand{\dotsscale}{1.8}

% Assets
\input{assets/geometry.tex}
\newcommand{\chapwithtoc}[1]{           %% Unumbered chapter included in TOC
  \chapter*{#1}
  \addcontentsline{toc}{chapter}{#1}
}

\newenvironment{solution}{
  \par\medskip\noindent
  \textit{Řešení}.
}{}

\newcommand{\cmark}{\ding{51}}
\newcommand{\xmark}{\ding{55}}

\MakeOuterQuote{"}

\newcommand{\topindent}{3mm}

\theoremstyle{definition}
\newtheorem{theorem}{Věta}[section]
\newtheorem{lemma}[theorem]{Lemma}
\newtheorem{proposition}[theorem]{Tvrzení}
\newtheorem{corollary}[theorem]{Důsledek}

\newtheorem{definition}[theorem]{Definice}
\newtheorem*{definition*}{Definice}
\newtheorem{remark}[theorem]{Poznámka}
\newtheorem*{remark*}{Poznámka}
\newtheorem*{symboldef}{Značení}
\newtheorem{example}[theorem]{Příklad}
\newtheorem*{example*}{Příklad}

% Zdrojové kódy
\definecolor{codebackground}{RGB}{247,248,250}
\definecolor{codecomment}{RGB}{63,127,95}
\definecolor{codekeyword}{RGB}{34,73,150}
\definecolor{codestring}{RGB}{150,58,58}
\lstdefinestyle{python}{
    language=Python,
    backgroundcolor=\color{codebackground},
    basicstyle=\ttfamily\small,
    keywordstyle=\color{codekeyword}\bfseries,
    commentstyle=\color{codecomment}\itshape,
    stringstyle=\color{codestring},
    numbers=left,
    numberstyle=\tiny\color{gray},
    numbersep=8pt,
    frame=single,
    rulecolor=\color{gray!45},
    framesep=5pt,
    breaklines=true,
    breakatwhitespace=false,
    showstringspaces=false,
    keepspaces=true,
    columns=fullflexible,
    tabsize=4,
    captionpos=t,
    upquote=true
}
\lstset{style=python}
\AtBeginDocument{%
    \renewcommand{\lstlistingname}{Program}%
    \renewcommand{\lstlistlistingname}{Seznam programů}%
    \counterwithin{lstlisting}{section}%
}

% Enumerate
\setlist[enumerate]{topsep=0pt,itemsep=-1ex,partopsep=1ex,parsep=1ex,label=\arabic*.}

% Itemize
\setlist[itemize]{topsep=0pt,itemsep=0ex,partopsep=1ex,parsep=1ex}

% Algoritmy
\numberwithin{algocf}{section}
\renewcommand{\algorithmcfname}{Algoritmus}
\SetKwInput{KwIn}{Vstup}
\SetKwInput{KwOut}{Výstup}
\SetKwProg{Function}{Funkce}{}{end}

% Barvy a zvýraznění převzaté z TI-prezentace/assets/macros.tex
\definecolor{lightblue}{HTML}{009AD4}
\definecolor{lightgreen}{HTML}{68FF00}
\definecolor{darkgreen}{HTML}{00B600}
\definecolor{darkblue}{HTML}{0265B3}
\definecolor{darkred}{HTML}{DA0707}
\definecolor{lightred}{HTML}{FF5100}
\definecolor{orange}{HTML}{FFE000}

\newcommand{\markred}[1]{\textcolor{lightred}{#1}}
\newcommand{\markdarkred}[1]{\textcolor{darkred}{#1}}
\newcommand{\markgreen}[1]{\textcolor{lightgreen}{#1}}
\newcommand{\markdarkgreen}[1]{\textcolor{darkgreen}{#1}}
\newcommand{\markorange}[1]{\textcolor{orange}{#1}}
\newcommand{\markblue}[1]{\textcolor{lightblue}{#1}}
\newcommand{\markdarkblue}[1]{\textcolor{darkblue}{#1}}

% Matematické značení sjednocené s TI-prezentace/assets/macros.tex
\newcommand{\R}{\mathbb{R}}
\newcommand{\C}{\mathbb{C}}
\newcommand{\N}{\mathbb{N}}
\newcommand{\Z}{\mathbb{Z}}
\newcommand{\Q}{\mathbb{Q}}

\newcommand{\T}[1]{#1^\top}
\newcommand{\compl}[1]{#1^\complement}
\newcommand{\mapping}[3]{#1:#2 \rightarrow #3}
\newcommand{\mapto}[3]{#1:#2 \mapsto #3}
\newcommand{\set}[1]{\left\{#1\right\}}
\newcommand{\powset}{\mathcal{P}}
\newcommand{\code}[1]{\langle#1\rangle}
\DeclareMathOperator*{\argmin}{arg\,min}
\DeclareMathOperator*{\argmax}{arg\,max}

\newcommand{\solutions}[1]{\left[K=\left\{#1\right\}\right]}
\newcommand{\sizeof}[1]{\left| #1 \right|}
\newcommand{\admid}{\;\middle\vert\;}
\newcommand{\abs}[1]{\left|\,#1\,\right|}
\newcommand{\norm}[1]{\left\lVert#1\right\rVert}
% Název \binary používá také balíček fmtcount, který může být v některých
% instalacích TeXu načten nepřímo. Specifičtější název zabrání kolizi.
\newcommand{\binarycode}[1]{\ensuremath{\left\langle #1\right\rangle}_B}

% Statistika a AI
\newcommand{\average}[1]{\overline{#1}}
\DeclareMathOperator{\median}{med}
\DeclareMathOperator{\Var}{var}
\DeclareMathOperator{\Cov}{cov}
\DeclareMathOperator{\MSE}{MSE}
\DeclareMathOperator{\SSE}{SSE}

\DeclareMathOperator{\TP}{TP}
\DeclareMathOperator{\TN}{TN}
\DeclareMathOperator{\FP}{FP}
\DeclareMathOperator{\FN}{FN}
\DeclareMathOperator{\Accuracy}{Acc}
\DeclareMathOperator{\Precision}{Prec}
\DeclareMathOperator{\Recall}{Rec}
\DeclareMathOperator{\Fscore}{F}
\DeclareMathOperator{\Gini}{Gini}
\DeclareMathOperator{\ReLU}{ReLU}

\DeclareMathOperator{\mspan}{span}
\DeclareMathOperator{\mrank}{rank}
\DeclareMathOperator{\tg}{tg}
\DeclareMathOperator{\id}{id}
\DeclareMathOperator{\dom}{dom}
\DeclareMathOperator{\rng}{rng}
\renewcommand{\emptyset}{\varnothing}

% Paths
\newcommand{\chapterspath}[1]{ch#1/content.tex}
\newcommand{\sectionspath}[1]{ch#1/sections}

% Info boxes
\newcommand{\infobox}[1]{\awesomebox[gray]{2pt}{\faInfoCircle}{gray}{\textcolor{gray}{#1}}}

% Seznamy používané v DFS a haldách
\DeclareMathOperator{\inlist}{in}
\DeclareMathOperator{\outlist}{out}
\DeclareMathOperator{\key}{key}

% A*
\newcommand{\Astar}{\ensuremath{\textup{A}^{*}}}

% Třídy složitosti a rozhodovací problémy
\DeclareMathOperator{\classP}{P}
\DeclareMathOperator{\classNP}{NP}
\DeclareMathOperator{\classNPC}{NPC}
\DeclareMathOperator{\classTime}{TIME}
\DeclareMathOperator{\classNTime}{NTIME}
\DeclareMathOperator{\classSpace}{SPACE}
\DeclareMathOperator{\classNSpace}{NSPACE}

\DeclareMathOperator{\SAT}{SAT}
\DeclareMathOperator{\THREESAT}{3-SAT}
\DeclareMathOperator{\UNSAT}{UNSAT}
\DeclareMathOperator{\NZMNA}{NzMna}
\DeclareMathOperator{\KLIKA}{Klika}
\DeclareMathOperator{\ROZVRH}{Rozvrh}
\DeclareMathOperator{\SUDOKU}{Sudoku}
\newcommand{\HALT}{\ensuremath{\textsc{Halt}}}
\newcommand{\FACTOR}{\ensuremath{\textsc{Factor}}}

\DeclareMathOperator{\regex}{RegE}

\setlength{\fboxsep}{3mm}
\newcommand{\problem}[3]{
  \fbox{
    \begin{minipage}{0.65\textwidth}
      \textsc{#1}\\[3mm]
        \indent\textbf{Vstup problému:} #2\\[3mm]
        \indent\textbf{Výstup problému:} #3
    \end{minipage}
  }
}

% Asymptotické odhady
\newcommand{\bigO}{\mathcal{O}}
\newcommand{\bigOmega}[1]{\Omega\left(#1\right)}

%hyperref
\hypersetup{
    colorlinks=true,
    urlcolor=black,
    bookmarksopen=true
}

% Parindent
\setlength{\parindent}{0em}
\setlength{\parskip}{1em}

% Logical operators
% \renewcommand{\implies}{\Rightarrow}
% \renewcommand{\iff}{\Leftrightarrow}

% Goniometrické funkce používané pouze ve skriptech
\DeclareMathOperator{\cotg}{cotg}

% Others
\newcommand{\rightnote}[1]{\hspace*{\fill} $\triangleleft$ \textit{#1}}
\newcommand{\degre}{\ensuremath{^\circ}}

% Images
\newcommand{\graphimgsize}{.75}

% TikZ makra převzatá z prezentací. V knižní třídě nepotřebujeme
% beamerovské parametry pro překryvy, názvy a vzhled však zůstávají stejné.
\tikzstyle{startstop} = [rectangle, rounded corners, minimum width=1.5cm,
  minimum height=0.5cm, text centered, draw=black, fill=red!30]
\tikzstyle{process} = [rectangle, minimum width=1.5cm,
  minimum height=0.5cm, text centered, draw=black, fill=orange!30]
\tikzstyle{decision} = [diamond, aspect=2, text centered,
  draw=black, fill=green!30]
\tikzstyle{arrow} = [thick, ->, >=stealth]

\tikzstyle{vertex} = [thick, draw, circle, minimum size=4mm,
  inner sep=1pt, outer sep=1pt, align=center]
\tikzstyle{vertexplain} = [minimum size=4mm, inner sep=1pt,
  outer sep=1pt, align=center]
\tikzstyle{verteximg} = [draw, circle, minimum size=4mm, inner sep=0pt,
  outer sep=1pt, align=center, clip]
\tikzstyle{verteximgplain} = [minimum size=4mm, inner sep=0pt,
  outer sep=1pt, align=center]
\tikzstyle{task} = [draw, rounded corners=2pt, minimum width=14mm,
  minimum height=8mm, inner sep=2pt, align=center]
\tikzstyle{edge} = [thick, >=stealth]
\tikzstyle{diredge} = [thick, ->, >=stealth]

\tikzstyle{state} = [draw, circle, thick, minimum size=10mm,
  inner sep=0pt, outer sep=0pt]
\tikzstyle{acceptstate} = [state, double, double distance=1pt]
\tikzstyle{transition} = [->, >=stealth, thick]

\newcommand{\state}[4][]{%
  \node[state,#1] (#2) at #3 {#4};%
}
\newcommand{\acceptingstate}[4][]{%
  \node[acceptstate,#1] (#2) at #3 {#4};%
}
\newcommand{\transitionarrow}[5][]{%
  \path (#2) edge[transition,#1] node[#3] {#4} (#5);%
}
\newcommand{\initialstate}[6][]{%
  \node[state,#1] (#2) at #3 {#4};%
  \draw[transition] #5 -- (#2.#6);%
}

% Chytrý křížový odkaz
\NewDocumentCommand{\smartref}{ O{} O{} m m m }{%
  \ifthenelse{\equal{#3}{true}}{%
    % Pokud je true, zobrazí se odkaz
    #1\ref{#4}#2% Zobrazí se prefix, odkaz a postfix
  }{%
    % Pokud není true, zobrazí se alternativní text
    #5%
  }%
}

\makeatletter
\renewcommand\@makefnmark{%
  \hbox{\@textsuperscript{\normalfont\@thefnmark}}%
}
\long\def\@makefntext#1{%
  \parindent 1em
  \noindent
  \hb@xt@1.8em{%
    \hss\@textsuperscript{\normalfont\ensuremath{\langle\@thefnmark\rangle}}%
  }#1%
}
\makeatother

\makeatletter
\@ifundefined{chapter}{}{%
  % Číslované kapitoly
  \renewcommand{\@makechapterhead}[1]{%
    \vspace*{50\p@}%
    {\parindent \z@ \centering
      \normalfont\large\scshape Kapitola \thechapter\par
      \vspace{1em}%
      \normalfont\huge\bfseries#1\par\nobreak
      \vskip 40\p@
    }}
  % Nečíslované kapitoly (beze změny popisku)
  \renewcommand{\@makeschapterhead}[1]{%
    \vspace*{50\p@}%
    {\parindent \z@ \centering
      \normalfont\huge\bfseries #1\par\nobreak
      \vskip 40\p@
    }}
}
\makeatother


% Title page info
\def\maintitle{Teoretická informatika}
\def\subtitle{Obor C,~3. ročník}
\def\authorname{David Weber}
\def\currentdate{\today}
\def\institution{SPŠE Ječná}
